\begin{summarybox}

	\textbf{\textcolor{CSummary}{一句话核心}}

	\textbf{离心率是连接焦点、准线和曲线形状的统一比例：几何定义为$e=\frac{|PF|}{d(P,l)}$；对椭圆和双曲线，还有代数形式$e=\frac{c}{a}$。}

	\medskip

	\textbf{\textcolor{CSummary}{使用条件}}

	\begin{itemize}[leftmargin=*, itemsep=0.5em, topsep=0.4em]

		\item \textbf{条件1（曲线类型）：} 适用于高中常见圆锥曲线：椭圆、双曲线、抛物线。

		\item \textbf{条件2（代数公式）：} $e=\frac{c}{a}$只适用于椭圆和双曲线；抛物线直接由定义得到$e=1$。

		\item \textbf{条件3（对应关系）：} 使用$e=\frac{|PF|}{d(P,l)}$时，焦点$F$和准线$l$必须对应。

	\end{itemize}

	\medskip

	\textbf{\textcolor{CSummary}{关键提醒}}

	\begin{itemize}[leftmargin=*, itemsep=0.5em, topsep=0.4em]

		\item \textbf{易错点（参数关系）：} 椭圆$c^{2}=a^{2}-b^{2}$，双曲线$c^{2}=a^{2}+b^{2}$，不可混淆。

		\item \textbf{易错点（准线方向）：} 焦点在$x$轴上，准线写成$x=\pm\frac{a^{2}}{c}$；焦点在$y$轴上，准线写成$y=\pm\frac{a^{2}}{c}$。

		\item \textbf{检查项（定义优先）：} 涉及焦点、准线、距离比例时，首先联想定义$e=\frac{|PF|}{d(P,l)}$。

	\end{itemize}

\end{summarybox}